\thispagestyle{empty}

\begin{center}
{\large\bfseries INDIGOScript \\ Aplicación web para el control del dispositivos astronómicos }\\
\end{center}
\begin{center}
Pablo Vallejo Ruiz\\
\end{center}

%\vspace{0.7cm}

\vspace{0.5cm}
\noindent\textbf{Palabras clave}: \textit{software libre, ingenieria de software, desarrollo web, TypeScript, Javascript, HTML, INDIGO, TCP/IP, patrones de diseño, gestión de proyectos, metodologias agile}
\vspace{0.7cm}

\noindent\textbf{Resumen}\\
La astronomía es un campo de estudio de notable complejidad que ha experimentado una evolución significativa a lo largo de las últimas décadas. Para poder satisfacer las exigencias del ámbito, y gracias a los avances tecnológicos actuales, los dispositivos de observación astronómica han experimentado un aumento exponencial en su sofisticación: hoy día ofrecen un mayor número de opciones para astrónomos tanto profesionales como aficionados, pero con ello su uso y gestión también se ha complicado considerablemente. 

Cámaras con dispositivos de carga acoplada (referidas como CCD, del inglés Charged-Couple Device), telescopios, motores de enfoque... Con el objetivo de automatizar y simplificar la gestión simultánea de toda clase de dispositivos distribuidos, se han desarrollado un sinfín de programas y protocolos que permiten el envío de mensajes entre estos. Sin embargo, en ocasiones estas soluciones no son compatibles con todos los dispositivos, no estan disponibles en su sistema operativo, o pueden ser demasiado complejas para algunos usuarios, que buscan algo más sencillo.

El presente proyecto tiene como objetivo ofrecer una solución simple y accesible que permita a los usuarios gestionar estos dispositivos de una manera intuitiva, a través del desarrollo de un cliente web que automatice la comunicación entre usuario y dispositivos. y una simple interfaz gráfica que haga uso de dicho cliente.


\cleardoublepage

\begin{center}
	{\large\bfseries INDIGOScript \\ Web app for the management of astronomical devices}\\
\end{center}
\begin{center}
	Pablo Vallejo Ruiz\\
\end{center}
\vspace{0.5cm}
\noindent\textbf{Keywords}: \textit{open source, software engineering, web development, TypeScript, JavaScript, HTML, INDIGO, TCP/IP, design patterns, project management, agile methodologies}, \textit{floss}
\vspace{0.7cm}

\noindent\textbf{Abstract}\\
Astronomy is a field of notable complexity that has undergone a significant evolution over the past few decades. To meet the demands of this field, and thanks to the current technological advancements, there has been an exponential increase in sophistication of astronomical obsrvation devices: nowadays, they offer more options than ever for both proffesional and amateur astronomers, but at the same time their use and managements has become more complex.

Charged-Couple Device (CCD) cameras, telescopes, mounts and focusers... In order to automate and manage all these kinds of distributed devices simultaneously, many dedicated programs and protocols have been developed. These solutions allow the exchange of control messages between user and devices. However, sometines they may be uncompatible with several devices, be unavailable in your operating system of choice, or just may be too complex for the necesities of more casual users.

The aim of this project is to provide a simple and accesible solution that allows users to manage these devices intuitively, by developing both a web client that automates communication between user and device, and a simple and approachable graphical interface that makes use of said web client.

\cleardoublepage

\thispagestyle{empty}

\noindent\rule[-1ex]{\textwidth}{2pt}\\[4.5ex]

D. \textbf{Sergio Alonso Burgos}, Profesor del Departamento de Lenguajes y Sistemas Informáticos de la Universidad de Granada 

\vspace{0.5cm}

\textbf{Informo:}

\vspace{0.5cm}

Que el presente trabajo, titulado \textit{\textbf{INDIGOScript}},
ha sido realizado bajo mi supervisión por \textbf{Pablo Vallejo Ruiz}, y autorizo la defensa de dicho trabajo ante el tribunal
que corresponda.

\vspace{0.5cm}

Y para que conste, expiden y firman el presente informe en Granada a Septiembre de 2024.

\vspace{1cm}

\textbf{El/la director(a)/es: }

\vspace{5cm}

\noindent \textbf{Sergio Alonso Burgos}

\chapter*{Agradecimientos}
%



